\begin{titlepage}
    \centering
    \vspace*{1.5cm}
    {\scshape\Large Universit\`a degli Studi --- Corso di Laurea Magistrale\par}
    \vspace{0.4cm}
    {\large Scalable and Reliable Systems\par}
    \vspace{2.0cm}
    \rule{\linewidth}{0.4pt}\\[0.6cm]
    {\huge\bfseries Intelligent Document Management System\par}
    \vspace{0.3cm}
    {\Large A Reliable, Multi-Tenant Platform with an MCP-Based Agentic Filing Layer\par}
    \rule{\linewidth}{0.4pt}\\[1.2cm]

    {\large\itshape Relazione tecnica di progetto\par}
    \vspace{0.4cm}
    {\large Dominio applicativo: \textbf{Finance \& Enterprise Document Compliance}\par}
    \vfill
    \begin{flushleft}
        \textbf{Repository}: \texttt{document-management-system}\\[4pt]
        \textbf{Data}: \today
    \end{flushleft}
\end{titlepage}

% ---------------------------------------------------------------------------
%  Abstract
% ---------------------------------------------------------------------------
\chapter*{Abstract}
\addcontentsline{toc}{chapter}{Abstract}

\noindent
Il presente lavoro descrive la progettazione, l'implementazione e la valutazione
di un \emph{Document Management System} (DMS) distribuito e multi-tenant, esteso
con uno strato di \emph{Agentic Operations} basato sul \emph{Model Context Protocol}
(MCP). Il sistema risolve un problema operativo concreto del dominio
\emph{Finance \& Enterprise Document Compliance}: l'acquisizione automatica,
la classificazione e l'archiviazione di documenti amministrativi (fatture,
contratti, ordini, comunicazioni di compliance) ricevuti via posta elettronica.

La trattazione adotta la struttura prevista dalle linee guida del corso, articolando
le motivazioni di dominio, l'architettura distribuita, i meccanismi di affidabilit\`a
e scalabilit\`a, l'osservabilit\`a, i ruoli dell'agente, il design degli MCP server,
la politica di sicurezza operativa, gli esperimenti di guasto, l'analisi
SLO--costi--\emph{ROI} e le politiche DevSecOps. Particolare attenzione \`e
riservata alla separazione fra automazione deterministica, decisione agentica
e \emph{human-in-the-loop}, alla gestione delle azioni ad alto impatto e
alla riduzione del rischio di allucinazione tramite \emph{structured outputs}
e soglie di confidenza.
